%! @@TITLE@@ — Unit @@UNITINT@@, Lesson @@LESSONID@@ — Lesson Plan
\documentclass[10pt]{article}
\usepackage{@@PREFIX@@-article}
\usepackage{@@PREFIX@@-boxes}
\usepackage{graphicx}   % -article does NOT load graphicx; needed for screenshots
\graphicspath{{images/}}
\newcommand{\TallMath}[1]{$\displaystyle #1\rule[-1.4em]{0pt}{3.2em}$}
@@COURSEMACROS@@\newcommand{\UnitNumberName}{Unit @@UNITINT@@: @@UNITTITLE@@ \quad}
\newcommand{\LessonNumberName}{Lesson @@LESSONID@@: @@TITLE@@}

\begin{document}
\begin{center}
    {\Large\bfseries \CourseName \\
    {\small\bfseries \UnitNumberName \LessonNumberName}}
\end{center}

\begin{tcolorbox}[colback=lilac, colframe=plum, boxrule=0.9pt, arc=2mm,
    left=2mm, right=2mm, top=1.5mm, bottom=1.5mm]
    \textbf{Primary Objective:} TODO --- restate the lesson aim in student terms.
    % Restate the CED Learning Objectives as student-facing aims (no "Big Idea" tag —
    % this CED has Mathematical Practices instead; see templates/lesson/course-workflow.md)
\end{tcolorbox}

% ── The clock, second thing the teacher reads ─────────────────────────────────
% GRADUAL RELEASE. The Min column MUST sum to exactly 60 (= \MeetingLength).
% Add it up by hand before you build — neither latexmk nor `make check` will.
% There is NO group-activity row: the whole release runs inside the guided notes.
% Model: unit01/lesson01/main.tex (the 2026-09-16 shape: no exit ticket, Homework Launch).
\boxguard[22]
\begin{skillbox}[Lesson Flow --- Gradual Release (60 minutes)]{lilac}
\small
\renewcommand{\arraystretch}{1.25}
\begin{tabularx}{\linewidth}{@{} l c X @{}}
\textbf{Phase} & \textbf{Min} & \textbf{What students are doing, and where it lives} \\
\midrule
Warm-Up                       & 5  & Silent, individual spiral review --- warm-up sheet. \\
Hook                          & 3  & Whole-class discussion --- hook box atop the guided notes. \\
\textbf{I Do} --- model       & 10 & Watch and annotate while the teacher thinks aloud --- notes \S TODO. \\
\textbf{We Do} --- guided     & 20 & Fill the notes together, every blank cold-called --- notes \S TODO. \\
\textbf{You Do} --- independent & 14 & Work alone while the teacher circulates --- notes \S TODO and the practice box. \\
Debrief                       & 5  & Whole-class share-out and the headline. \\
Homework Launch               & 3  & Start homework problems TODO alone; the teacher reads the diagnostic item over shoulders. \\
\end{tabularx}

\smallskip
\textbf{If the clock slips:} TODO --- name what to cut FIRST and what to protect.
Never pad the total to reach 60; cut, and say here what you cut.
\end{skillbox}

\renewcommand{\arraystretch}{1.6}
\begin{skillbox}[Priority Ideas \& Skills]{goldbox}
    % TODO: left column = priority skills; right column = Key Understandings
    %       (paraphrased from the EKs).
\end{skillbox}

\begin{skillbox}[{Vocabulary, Concepts, \& Theorems}]{greenbox}
    % TODO: key terms + definitions for this lesson.
\end{skillbox}

\begin{skillbox}[Activate Prior Knowledge \& Spiral Review (5 min)]{lilac}
    Please see attached warm-up (spiral review).
    % TODO: list the prerequisite skills this warm-up reviews, in words.
    % Text only — never embed a warm-up thumbnail (\includegraphics{warmup/main}),
    % whether the warm-up is authored or prefab. See LESSON_SHAPE.md §4.
\end{skillbox}

\begin{skillbox}[Hook (3 min)]{lilac}
    % TODO: the engaging entry question / scenario.
\end{skillbox}

% ── The lesson, tagged by release phase ───────────────────────────────────────
% Every part opens with its phase in CAPS, the part title, and its minutes, then
% an italic line naming the notes section it covers AND who is holding the pen.
% A part that does not say who is holding the pen is not tagged.
\begin{skillbox}[Lesson --- I Do / We Do / You Do]{lilac}
\begin{multicols}{2}

\textbf{I DO --- Part 1: TODO} (10 min)

\textit{Guided notes \S TODO. Teacher works; students watch and annotate. Nothing
is cold-called in this phase.}
\begin{itemize}
  \item % TODO
\end{itemize}

\textbf{WE DO --- Part 2: TODO} (10 min)

\textit{Guided notes \S TODO. Class fills the blanks together; cold-call every one.}
\begin{itemize}
  \item % TODO
\end{itemize}

\columnbreak

\textbf{WE DO --- Part 3: TODO} (8 min)

\textit{Guided notes \S TODO. Still together, but hand over more of each answer.}
\begin{itemize}
  \item % TODO
\end{itemize}

\textbf{YOU DO --- Part 4 and Practice: TODO} (14 min)

\textit{Guided notes \S TODO and the practice box. Students work alone; circulate
and say nothing a neighbor could say instead.}
\begin{itemize}
  \item % TODO
\end{itemize}
\end{multicols}
\end{skillbox}

\begin{skillbox}[Active Monitoring]{redbox}
Circulate during the \textbf{You Do} phase --- the stretch where students are
working without you --- and check:
    % TODO: common errors + cold-call prompts.
\end{skillbox}

% ── Differentiation — INSIDE the You Do phase. There is no activity/ component. ─
\begin{skillbox}[Differentiation --- During You Do (14 min)]{redbox}
    % TODO: Support / On level / Extend. These are moves the teacher makes while
    % circulating the practice box, NOT a separate handout: which item to sit
    % with a struggling student on, which to let the class run, and which
    % question to hand a finisher.
\end{skillbox}

% ── Debrief — TEACHER-FACING ONLY. There is no debrief/ student component. ─────
\boxguard[20]
\begin{skillbox}[Debrief (5 min)]{redbox}
Whole class, teacher-led, packets still open. Three moves, in order:
\begin{enumerate}[itemsep=3pt]
  \item \textbf{Share out the You Do} (2 min) --- one answer per practice item,
        not a full review, so a student who never reached the last item still
        hears its result.
        % TODO: name the one answer you want from each practice item.
  \item \textbf{Name the headline} (1 min) --- say it, then make the class say it
        back. % TODO: the day's central sentence, verbatim, in italics.
  \item \textbf{Point forward} (2 min) --- % TODO: next lesson, or a read-only
        notes section read aloud here rather than during notes.
\end{enumerate}
\textbf{Do not} re-teach here, take new questions, or let the homework start
early. A debrief that runs long eats the Homework Launch, the period's last formative read.
\end{skillbox}

\begin{skillbox}[Homework Launch (3 min)]{redbox}
\textbf{There is no exit ticket.} Due date on the board and into the cover's Homework
row; students start homework problems % TODO: which problems
alone and in silence while you circulate.
    % TODO: name the DIAGNOSTIC homework item (the day's central distinction) and
    % what a miss looks like; then the three piles to sort what you see into.
\end{skillbox}

\begin{skillbox}[Reinforcement \& Extension]{goldbox}
\textbf{Homework --- printed, graded (2 pages).} % TODO: overview by part and
% problem number; every context new; name the diagnostic item and the headline item.

\textbf{AP Practice --- extra credit (2 pages).} % TODO: Section I (five MC, four
% options) --- what each distractor catches; Section II (one FRQ) --- what each part asks.

\textbf{Preview:} % TODO: next lesson.
\end{skillbox}

% ── Teacher notes — one per component, in packet order ────────────────────────
% Teacher-only prose lives HERE, never in a _key: a note is the one block in a
% key with no counterpart in the blank, so it makes the key longer than the
% blank. See the shared skill's references/conventions.md ("teachernote").
\begin{teachernote}[Warm-Up]
    % TODO: pacing, common errors, what to look for.
\end{teachernote}

\begin{teachernote}[Guided Notes]
    % TODO: which phase to slow down for, and what to cut if the clock slips.
\end{teachernote}

\begin{teachernote}[Debrief]
    % TODO: why these minutes are worth protecting, and what to borrow from
    % instead when the guided phase overruns (the last practice item, not this).
\end{teachernote}

\begin{teachernote}[AP Practice]
    % TODO: extra credit, optional; the MC answer string; the distractor worth
    % looking for; how to credit the FRQ's stretch part.
\end{teachernote}

\begin{teachernote}[Homework]
    % TODO: which items to grade for accuracy, which predict the next lesson, and
    % the common slip on each and what to do about it in the next warm-up.
\end{teachernote}
\end{document}
